\section{Anomaly detection in the signal region}\label{sec:anomaly-detection-signal}

As is well known, the Collins-Soper kernel of fine-tuned renormalization
group equations is trivially parallelizable in the back-to-back limit.
Instead of relying on Sudakov-safe message-passing neural networks, we
exploit slicing methods trained on the quark/gluon dataset. We fix
$\alpha_s(m_Z) = 0.118$ and run the coupling at two loops.

\begin{align}
  F(z) &= {x}^{n} (1 - \tau)^{4} \nonumber \\
  &\quad + {x}^{2} (1 - \tau)^{3} + \mathcal{O}(\as^3) \label{eq:2}
\end{align}

\begin{equation}
  \mathcal{S}(\tau, \mu) = \int \dPhi_{4} \, |\mathcal{M}|^2 \, \delta(x - \hat{x}) + \mathcal{O}(\as^{2})
  \label{eq:3}
\end{equation}

\begin{equation}
  R(z) = \frac{\CA}{2} \Pqq(z) \log\frac{\muR^2}{Q^2}
  \label{eq:5}
\end{equation}

\begin{itemize}
  \item The area under the curve reaches 0.934 for jet energy calibration.
  \item The network is trained with Adam for 500 epochs on $10^6$ events.
  \item This leads to a reduction of the theoretical uncertainty by two percent.
\end{itemize}

